\documentclass[11pt,a4paper]{article}
\newcommand{\labonaamwaarde}{labo 4}
\usepackage{syllabus}

\begin{document}

\syllabustitel{labo 4 -- \emph{Dataverwerking II}}{Combineren, samenvatten, herschikken}

% ---------------------------------------------------------------------------
% Gedeeld door de syllabi van labo 3, 4 en 5. Pas dit hier aan, niet in de labobestanden.
\section{Overzicht}

Tijdens labo 3, 4 en 5 ben jij de nieuwe data-analist van een
(fictieve) jeugdclub voor voetbal en atletiek. De hoofdtrainer bezorgt je
drie exportbestanden uit de systemen van de club. Die zijn bewust
rommelig, zoals echte exports: inconsistente labels, sprinttijden in
milliseconden, ontbrekende waarden, een atleet die twee keer voorkomt en
een kolom met een komma als decimaalteken.

Elke sessie bouwt verder op de vorige:

\begin{description}[leftmargin=1.6cm, style=nextline]
  \item[Labo 3] \textbf{Van bestand tot nette tabel.} Inladen, verkennen,
    filteren, kolommen toevoegen en opkuisen tot je de data kunt
    vertrouwen.
  \item[Labo 4] \textbf{Combineren, samenvatten, herschikken.} Groepen
    samenvatten, tabellen koppelen, wisselen tussen een brede en een lange
    vorm, en werken met een hartslagreeks per seconde.
  \item[Labo 5] \textbf{Vertel het verhaal met plotly.} De juiste grafiek
    kiezen voor een vraag, eerlijke grafieken maken met
    \code{plotly.express} en ze combineren tot een klein interactief
    dashboard.
\end{description}


% ---------------------------------------------------------------------------
\section{Leerdoelen}

Na labo 4 kun je:

\begin{enumerate}[label=\textbf{LD\arabic*}, leftmargin=1.4cm]
  \item data per groep samenvatten met \code{groupby} en \code{agg}, en
        tabellen koppelen met \code{merge} zonder dubbel te tellen;
  \item data herschikken tussen een lange en een brede vorm met
        \code{melt} en \code{pivot\_table};
  \item eenvoudige kenmerken van een tijdreeks berekenen: hartslagzones,
        voortschrijdende gemiddelden en hersampling.
\end{enumerate}

% ---------------------------------------------------------------------------
\section{Voorkennis}

Labo 3 (\emph{Dataverwerking I}): een CSV-bestand inladen met
\code{pandas}, rijen en kolommen selecteren en filteren, nieuwe kolommen
berekenen en de data opkuisen. Labo 4 start met de opkuis-pipeline van
labo 3, in compacte vorm.

% ---------------------------------------------------------------------------
% Gedeeld door de syllabi van labo 3, 4 en 5. Pas dit hier aan, niet in de labobestanden.
\section{Software en lesmateriaal}

Elk labo is een \textbf{marimo-notebook} in de map \code{labs/} van de
cursusrepository. Je kunt het op twee manieren openen:

\begin{itemize}
    \item \textbf{Online:} surf naar \url{https://marimo.app}, kies
        \emph{New $\rightarrow$ Open from URL\dots} en plak de raw-URL van
        het notebook op GitHub. De data wordt dan van GitHub gelezen, dus
        je hoeft niets te installeren.
    \item \textbf{Lokaal:} installeer de pakketten uit
        \code{requirements.txt} en open het notebook met marimo. De data
        wordt dan uit de map \code{data/} gelezen.

\end{itemize}

Drie dingen die je moet weten over marimo:

\begin{itemize}
  \item Notebooks zijn \textbf{reactief}: als je een cel aanpast, worden
        alle cellen die ervan afhangen opnieuw uitgevoerd, in de volgorde
        van hun afhankelijkheden.
  \item \textbf{Twee cellen kunnen niet dezelfde variabele definiëren.}
        Geef elke stap een eigen, duidelijke naam (\code{tests\_raw},
        \code{tests\_clean}, \dots) of zet een reeks stappen samen in één
        cel.
  \item Elke oefening bevat een placeholder \code{...}. Een verborgen cel
        eronder controleert je antwoord: die kleurt \textbf{groen} als het
        juist is, en geeft anders een algemene hint over wat nog niet
        klopt, zonder het antwoord te verklappen.
\end{itemize}


% ---------------------------------------------------------------------------
% Gedeeld door de syllabi van labo 3, 4 en 5. Pas dit hier aan, niet in de labobestanden.
\section{De dataset}

\noindent
\rowcolors{2}{tint}{white}
\begin{tabularx}{\textwidth}{@{}L{3.1cm}L{3.2cm}Y@{}}
\toprule
\textbf{Bestand} & \textbf{Omvang} & \textbf{Inhoud en gekende problemen} \\
\midrule
\code{athletes.csv}   & 43 rijen, 7 kolommen   & ID, naam, geboortedatum, geslacht, positie, lengte, gewicht. 13 schrijfwijzen voor 4 posities, gewicht als tekst met een decimale komma (\code{"57,5"}), één atleet twee keer ingevoerd met twee ID's, enkele ontbrekende waarden. \\
\code{tests.csv}      & 114 rijen, 8 kolommen  & Sprint over 10\,m en 30\,m, spronghoogte (countermovement jump) en Yo-Yo-niveau, op 3 testmomenten. Sommige sprinttijden staan in milliseconden; enkele ontbrekende waarden. \\
\code{hr\_session.csv} & 12\,000 rijen, 5 kolommen & Hartslag en snelheid per seconde van 5 atleten tijdens één training van 40 minuten. \\
\bottomrule
\end{tabularx}
\rowcolors{1}{}{}


% ---------------------------------------------------------------------------
\needspace{18\baselineskip}
\section{Planning}

% Planning van één labo. Gebruikt door de syllabus van dit labo en door de syllabus voor labo 3-5.
\textbf{Doel:} data aggregeren met \code{groupby}, tabellen koppelen met
\code{merge}, herschikken met \code{melt} en \code{pivot}, en eenvoudige
kenmerken van een tijdreeks berekenen.

\begin{schedule}
13:00--13:15 & Herhaling: een kapotte cel herstellen & Een kleine opkuis-pipeline met twee fouten repareren. \\
13:15--14:10 & \code{groupby} + \code{agg} & Statistieken per groep, leeftijdscategorieën met \code{pd.cut}, meerdere benoemde aggregaties tegelijk. \\
14:10--14:55 & \code{merge} & Inner en left joins; waarom een dubbele atleet twee keer meetelt, en waarom je dubbels verwijdert vóór het koppelen. \\
14:55--15:10 & Pauze & \\
15:10--15:50 & \code{melt} en \code{pivot} & Lange en brede vorm; vooruitgang tussen twee testmomenten. \\
15:50--16:50 & Tijdreeksen & Hartslagzones als \% van $HR_{max}$ (220 $-$ leeftijd), tijd per zone, \code{rolling}, \code{resample}. \\
16:50--17:30 & Minicase per twee & Een trainingsrapport: één rij per atleet. \\
\end{schedule}

\textbf{Oefeningen:} (0) herhaling: de kapotte cel herstellen;
(1) groeperen per geslacht en leeftijdscategorie; (2) wat heeft het
verwijderen van de dubbels gekost?; (3) vooruitgang in spronghoogte;
(4) minuten in zone 4--5 per atleet; \textbf{minicase:}
trainingsrapport met de minuten in zone 4--5, de maximale snelheid en de
afgelegde afstand per atleet.


% ---------------------------------------------------------------------------
\section{Werkwijze}

\begin{itemize}
  \item Elk blok start met een korte demo in het notebook, gevolgd door
        oefeningen die je in het notebook zelf oplost.
  \item Gebruik de automatische controles als feedback, niet als doel:
        als een controle groen kleurt, zorg dan dat je kunt uitleggen
        \emph{waarom} je code werkt.
  \item De minicase op het einde maak je per twee.
\end{itemize}

% ---------------------------------------------------------------------------
\section{Voorbereiding}

\begin{itemize}
  \item Controleer vóór de sessie dat je het notebook kunt openen,
        lokaal of op marimo.app.
  \item Zorg dat je elke stap van de opkuis-pipeline op het einde van
        labo 3 kunt uitleggen, want labo 4 vertrekt daarvan.
  \item Lees de cursustekst bij labo 4 naast het notebook.
\end{itemize}

\end{document}
